Dans un temps économique marqué par une concurrence plus forte et une demande qui monte pour la qualité et le bon travail, les entreprises qui font des choses cherchent tout le temps à améliorer leurs façons de faire et à s'assurer que leurs machines sont là quand il faut. La bonne gestion de la réparation et la bonne manière de gérer la production sont devenues des gros problèmes pour garder la compétitivité et la vie longue des groupes.

\vspace{1cm}

De nos jours, les systèmes d’information ont un rôle important dans cette façon d’améliorer sans cesse. L’automatisme et la numérisation des procédés aident non seulement à augmenter la productivité, mais aussi à garantir une meilleure trace et une décision plus rapide et plus sûre. C’est dans ce contexte que s'inscrit ce projet de fin d'études, qui est de faire une application web de maintien et de production.
\vspace{1cm}


Le but de ce projet est de faire et de mettre en place une plate-forme web pour gérer les travaux d'entretien (prévus et imprévus) ainsi que le suivi des actions produites. Cette app aidera à planifier les interventions, suivre les machines, gérer des ordres de fabrication et voir les résultats grâce à des tableaux de bord faciles à comprendre.

\vspace{1cm}


Pour frapper ces buts, le travail utilise des technologies du web nouvelles qui garantissent performance, accès et sécurité. Il parle surtout aux techniciens de maintenance, responsables de production et chefs qui voudraient avoir un outil commun pour voir en direct leurs tâches industrielles.

\vspace{1cm}
